\chapter{Introduction}
\label{ch:intro}

\section{Background}
\label{sec:intro_background}

Reconstructing a continuous turbulent flow field from sparse point measurements arises in aerospace boundary-layer estimation, wind-farm wake monitoring, biomedical Particle Image Velocimetry (PIV) slice fusion, and active flow control. Field deployment in these settings shares three operational constraints. First, instruments report at most $\sim 10^2$ point measurements, whereas the field to be reconstructed carries far more degrees of freedom (of order $10^5$ on the $256^2$ grid used here, growing as $Re^{9/4}$ for a three-dimensional simulation at high Reynolds number~\cite{Pope2000}). Second, no offline DNS reference is available at training or inference. Third, low-latency inference from the live sensor feed is required for active flow control.

Approaches to estimating a flow field from limited measurements have evolved through three broad stages. Classical reduced-order models compress the field onto a handful of modes identified offline and place sensors to condition that basis, while variational and ensemble data assimilation~\cite{Asch2016DA} fuse measurements with a forward solver; both rely on either an offline trajectory or repeated solver calls. Supervised deep-learning surrogates~\cite{VinuesaBrunton2022NCS} learn the measurement-to-field map directly in a single forward pass, but require paired full-field simulation data. Physics-informed methods relax that data demand: physics-informed neural networks~\cite{Raissi2019JCP} embed the PDE residual as a training constraint, neural operators~\cite{Lu2021NMI} learn mesh-free solution operators queryable at any coordinate, and physics-informed neural operators~\cite{Li2021PINO} combine the two, yet these target mostly forward or dense-input problems, not the sparse-sensor inverse setting. No line jointly satisfies the field-deployable constraints above (sensor-only training without a full reference field, and low-latency inference from a sparse live feed), which is the regime this work targets; \S\ref{sec:literature} examines each line in turn.

The difficulty underlying every stage has a quantitative anchor: reconstructing $N$ field values from $K \ll N$ point measurements is an under-determined inverse problem, and compressed-sensing theory bounds when exact recovery is achievable. The compressed-sensing bound~\cite{Donoho2006, Candes2006},
\begin{equation}
    \label{eq:cs_bound}
    M \;\ge\; C_0 \, s \, \log_2(N/s),\qquad C_0\in[1,2],
\end{equation}
where $M$ is the number of measurements, $s$ the number of significant wavelet coefficients, and $N$ the grid size, places the exact-recovery threshold for sparse signals on a $256^2$ grid an order of magnitude above the typical $\sim 10^2$ sensor budget. Engineering deployment at $K \sim 10^2$ therefore does not seek pointwise exact recovery; it asks how much of the energy-dominant flow structure can be reconstructed within this budget under a low-latency inference constraint, and whether the residual error is bounded by the sensor budget or by the chosen architecture.\label{sec:underdetermined}

\section{Literature Review}
\label{sec:literature}

This section reviews seven research lines on sparse-sensor turbulent-flow reconstruction, ordered from classical to recent and from algorithm-only to physics-aware (Table~\ref{tab:lit_summary}). \S\ref{sec:gap} consolidates the methodological limitations into four research gaps.

\begin{table}[!htbp]
    \centering
    \caption{Seven research lines relevant to sparse-sensor turbulent-flow reconstruction and the structural limitation that blocks each from addressing the engineering-deployable setting on its own.}
    \label{tab:lit_summary}
    \footnotesize
    \setlength{\tabcolsep}{3pt}
    \renewcommand{\arraystretch}{1.1}
    \begin{tabularx}{\textwidth}{
        >{\raggedright\arraybackslash}p{2.0cm}
        >{\raggedright\arraybackslash}p{3.4cm}
        >{\raggedright\arraybackslash}X
        >{\raggedright\arraybackslash}p{1.9cm}
        >{\raggedright\arraybackslash}p{2.3cm}
    }
    \toprule
    \textbf{Research line} & \textbf{Representative works} & \textbf{Main contribution} & \textbf{Supervision regime} & \textbf{Structural limitation} \\
    \midrule
    Reduced-Order Model (ROM) \& sparse identification
    & POD~\cite{Sirovich1987POD}; DMD~\cite{Schmid2010DMD}; SINDy~\cite{Brunton2016SINDy}; QR-pivot~\cite{Manohar2018PNAS}
    & Compress the field onto a low-rank basis identified offline.
    & Offline DNS trajectory
    & Needs offline trajectory; linear basis. \\
    \midrule
    Data assimilation
    & 4D-Var, EnKF~\cite{Asch2016DA}; B-PINN~\cite{Yang2021JCP}
    & Bayesian state estimation against a forward solver and noise model.
    & Forward solver + sensor
    & Adjoint cost; HMC scaling at high $Re$. \\
    \midrule
    Deep super-resolution / ROM
    & Fukami, 2019~\cite{Fukami2019JFM}; Maulik, 2021~\cite{Maulik2021NN}
    & Learn an upsampling map or latent dynamics from paired snapshots.
    & Paired DNS supervision
    & Paired DNS required; no PDE consistency. \\
    \midrule
    Operator learning
    & DeepONet~\cite{Lu2021NMI}; FNO~\cite{Li2021FNO}; Galerkin/OFormer~\cite{Cao2021Galerkin, Li2023OFormer}
    & Learn a mapping between whole functions (e.g.\ forcing $\to$ solution field): a branch sub-network reads the input function, a trunk evaluates the output at any query coordinate.
    & Paired field supervision
    & Dense-input forward operators; not sparse inverse. \\
    \midrule
    Stabilized PINNs
    & Wang 2022~\cite{wang2022respectingcausality}, 2024~\cite{Wang2024JMLR}, 2025~\cite{wang2025gradientalignmentpinns}
    & Causal weighting, PirateNet, and SOAP make stiff high-$Re$ training tractable.
    & PDE-only (no sensor)
    & Forward focus; MLP backbone underuses the sensor trajectory. \\
    \midrule
    Liquid NN / continuous-time
    & LTC~\cite{Hasani2021LTC}; CfC~\cite{Hasani2022CfC}; Neural/Latent ODE~\cite{Chen2018NODE, Rubanova2019LatentODE}
    & ODE-based and closed-form RNN dynamics for irregular sampling.
    & Paired sequence supervision (no PDE)
    & Not applied to PDE or fluid reconstruction. \\
    \midrule
    Sparse-sensor with physics
    & Mo \& Magri, 2025~\cite{Mo2025PRF}; Kelshaw, 2022~\cite{Kelshaw2022PISR}; Parfenyev, 2024~\cite{Parfenyev2024PINN}; SHRED~\cite{Williams2024SHRED}; Senseiver~\cite{Senseiver2023}
    & Most recent sensor-only-with-PDE or DNS-paired architectures; the first three use sensor measurements plus the PDE only.
    & Sensor-only with PDE (Mo \& Magri, Kelshaw, Parfenyev) vs.\ DNS-paired with priors (others)
    & Sensor-only-with-PDE works are grid-based at $Re\!\lesssim\!10^3$; no query-anywhere continuous-time operator variant at high $Re$. \\
    \bottomrule
    \end{tabularx}
\end{table}
\FloatBarrier

\paragraph{Classical reduced-order and data-assimilation methods.}
Proper Orthogonal Decomposition (POD)~\cite{Sirovich1987POD}, Dynamic Mode Decomposition (DMD)~\cite{Schmid2010DMD}, and sparse identification~\cite{Brunton2016SINDy} compress the field onto a low-rank linear basis, with QR-pivot~\cite{Manohar2018PNAS} or Q-DEIM~\cite{Drmac2016SISC} placement conditioning that basis; all require an offline DNS trajectory whose cost grows steeply with Reynolds number. Variational and ensemble Bayesian assimilation~\cite{Asch2016DA} instead fuse measurements with a forward solver, trading trajectory cost for repeated adjoint or sampling calls. Neither delivers the fast, reference-free inference a live control loop demands.

\paragraph{Learned reconstruction trained on full fields.}
A second line learns the measurement-to-field map directly. Convolutional super-resolution~\cite{Fukami2019JFM} and recurrent-autoencoder reduced-order models~\cite{Maulik2021NN} upsample or evolve a latent from paired snapshots; more recent sensor-input designs improve accuracy: SHRED~\cite{Williams2024SHRED} regresses the full state from a sensor time series through a shallow decoder, Senseiver~\cite{Senseiver2023} and the Energy Transformer~\cite{EnergyTransformer2025} use cross-attention or associative-memory inpainting, and FLRNet~\cite{FLRNet2024} learns a Fourier-feature variational autoencoder. Whatever the accuracy, every one of these is supervised by a full reference field (as a decoder target, a perceptual or variational loss, or masked-patch inpainting) and imposes no PDE, so none can be trained where no DNS field exists (Gap~1).

\paragraph{Operator learning for partial differential equations.}
Operator networks learn mappings between whole functions. DeepONet~\cite{Lu2021NMI} factorises the operator into a branch that reads the input function and a trunk that evaluates any query coordinate~\cite{ChenChen1995} (the query-anywhere readout adopted in this work), while the Fourier Neural Operator (FNO)~\cite{Li2021FNO} and attention variants~\cite{Cao2021Galerkin, Li2023OFormer} apply global spectral or kernel mixing. These are demonstrated as dense-input forward operators: the branch expects its input function sampled on a fixed dense grid (the FNO spectral mixing requires a regular mesh) rather than $\sim 10^2$ scattered points, and supervision is paired input--output fields. Even the physics-informed operator PINO~\cite{Li2021PINO} evaluates its residual on a grid. The sparse, scattered-sensor inverse problem lies outside this regime (Gap~2).

\paragraph{Physics-informed and continuous-time networks.}
Physics-Informed Neural Networks (PINNs)~\cite{Raissi2019JCP} impose the PDE residual as a soft constraint and need no full field, and Hidden Fluid Mechanics~\cite{Raissi2020Science} extends them to indirect observation. High-$Re$ stiffness from loss-gradient imbalance~\cite{Krishnapriyan2021} is eased by causal weighting~\cite{wang2022respectingcausality}, PirateNet~\cite{Wang2024JMLR}, and SOAP preconditioning~\cite{wang2025gradientalignmentpinns}, which have carried PINN training to $Re\sim10^4$. Two obstacles persist for the present setting: the coordinate-Multi-Layer-Perceptron (MLP) backbone takes only $(\mathbf{x},t)$ and sees the sensors solely through a loss term, never reading the measurement stream as an input (Gap~2); and causal or curriculum schemes optimise sequentially in time and per instance, so each new flow requires a full re-training with no amortised query. Continuous-time recurrent cells, the Closed-form Continuous-time (CfC) variant of liquid networks~\cite{Hasani2022CfC}, provide input-driven dynamics that tolerate irregular sampling at single-step cost, unlike solver-based Neural Ordinary Differential Equations (Neural ODEs)~\cite{Chen2018NODE}, but have served control and time-series tasks rather than PDE-constrained spatial reconstruction (Gap~3).

\paragraph{Closest prior work, and the remaining gap.}
Three recent works share the field-reproducible supervision adopted here: sensor measurements and the PDE residual, with no full reference field. Mo and Magri~\cite{Mo2025PRF} reconstruct a turbulent Kolmogorov flow with a physics-constrained convolutional network; Kelshaw, Rigas and Magri~\cite{Kelshaw2022PISR} super-resolve sparse Kolmogorov observations with a physics-informed CNN; and Parfenyev, Blumenau and Nikitin~\cite{Parfenyev2024PINN} reconstruct forced two-dimensional turbulence with a PINN. All three show that sensor-only-with-physics supervision is viable, but each works on a grid and at a far lower Reynolds number than the present $Re=10^4$ ($Re\approx34$ to $\approx10^3$): they return a fixed mesh rather than a query-anywhere field, and form the residual by finite differences or grid-restricted operators rather than continuous automatic differentiation. The query-anywhere architectures conversely require dense full-field supervision: the Senseiver~\cite{Senseiver2023}, and the spatio-temporal deep operator network of Vo Dang and Nguyen~\cite{FLRONet2024} (the closest published architecture to the present branch--trunk readout) nonetheless train against complete CFD fields rather than the PDE. Among the surveyed methods, none combines query-anywhere continuous-time evaluation with sensor-only-with-physics training at this Reynolds number, which is the cell this work occupies. Reported gains in this area are also sensitive to baseline strength~\cite{McGreivy2024}, motivating the fair-baseline and multi-seed protocol used throughout.

\section{Motivation}
\label{sec:gap}

The reconstruction question raised in \S\ref{sec:intro_background} (how much of the flow $K\sim 10^2$ sensors can recover, and whether the residual is bounded by the sensor budget or by the chosen architecture) cannot be settled with any surveyed method, because each leaves at least one of four obstacles open. The stakes are practical: closed-loop flow control and live wake monitoring need a full-field estimate in real time from the few sensors a rig carries, yet the surveyed methods force a choice between a DNS field unavailable in deployment and interpolation or forward solvers too slow or too inaccurate for the control loop. The first three gaps are engineering obstacles the reconstruction architecture must close; the fourth is the open question of how the sensing configuration (sensor count, placement, and noise) governs the achievable result. Each gap is stated below with the response the proposed work provides: a new architecture for Gaps 1--3, and a systematic sensing study for Gap 4.

\paragraph{Gap 1: No ground-truth field to train against.}
On a real rig there is no full-field reference: only the sensor readings and the governing equations are available. Most surveyed methods train against a DNS full field through perceptual, latent, or VAE losses (learned image-similarity or generative-model objectives~\cite{Williams2024SHRED, Senseiver2023, FLRNet2024, EnergyTransformer2025}), which cannot be reproduced in field deployment; only Mo and Magri~\cite{Mo2025PRF} restrict supervision to sensor measurements plus the PDE at both training and inference. The proposed framework targets this field-reproducible regime at the $K \sim 10^2$ sensor budget typical of industrial Particle Image Velocimetry (PIV) and hot-wire installations.

\paragraph{Gap 2: The model must read the sensor stream, not merely be scored by it.}
With no full field to imitate, the network must turn the handful of point readings into the whole field, so the measurements must enter as an \emph{input the model consumes}, not as a penalty term. A standard Physics-Informed Neural Network (PINN) backbone~\cite{wang2025gradientalignmentpinns,Wang2024JMLR} sees the sensors only on the loss side and never reads them as input; operator networks~\cite{Lu2021NMI,Li2021FNO,Kovachki2023JMLR} do accept a function-valued input but expect it sampled densely on a fixed grid, not $\sim 10^2$ scattered points. No surveyed backbone ingests sparse, scattered sensors as its input.

\paragraph{Gap 3: Sensors report on uneven clocks, yet the model must stay differentiable for the PDE.}
Field sensors report at irregular, unsynchronised time stamps, so a deployable temporal model must tolerate arbitrary gaps, a capability retained here even though the controlled Kolmogorov benchmark samples uniformly. The model must also be twice-differentiable in space and time, because evaluating the Navier--Stokes residual requires second spatial derivatives of the network output (the viscous term $\nu\nabla^2\mathbf{u}$), computed via automatic differentiation through the network. ODE-based temporal bridges~\cite{Chen2018NODE,Rubanova2019LatentODE} become unstable under second-order differentiation, whereas the Closed-form Continuous-time (CfC) cell~\cite{Hasani2022CfC} preserves continuous-time stability at the cost of a single recurrent step but has not been applied to PDE-constrained reconstruction. Together with Gap 2, this motivates a hybrid backbone combining a DeepONet branch--trunk factorisation, CfC recurrence, and a distance-aware cross-attention readout.

\paragraph{Gap 4: How the sensing configuration (count, placement, and noise) governs reconstruction is unmapped.}
Surveyed methods report a single aggregate error for one fixed sensing setup, neither separating what the sensors can carry from what the model extracts nor tracing how the result changes with the configuration. Three levers remain unquantified in the final result set: sensor \emph{count}, which sets an information-theoretic ceiling on resolvable wavenumbers; sensor \emph{placement}, which determines how much of that information is captured; and measurement \emph{noise}, which tests whether the reconstruction is fragile to realistic observation perturbations. Without this map, it is not possible to determine whether more sensors, better placement, cleaner sensing, or a better model is the productive next step. The proposed work treats the sensing configuration as a controlled variable and anchors the count axis to a per-wavenumber-band ceiling from the compressed-sensing threshold~\eqref{eq:cs_bound} and measured db4 wavelet statistics, so each residual error can be attributed to a sensing limit or to architectural headroom.

\section{Objective}
\label{sec:research_questions}

This work introduces PI-CON (Physics-Informed Continuous-time Operator Network), a CfC--DeepONet hybrid that reconstructs an unsteady 2D turbulent velocity field from sparse point sensors and the Navier--Stokes equations alone (trained with no DNS full field and queried in real time), then maps how the sensing configuration governs that reconstruction. PI-CON replaces the offline-DNS trajectory of POD-based reduced-order models and the iterative adjoint-and-solver loop of 4D-variational assimilation with a single learned forward pass over the live sensor stream. Three objectives pair a qualitative goal with a falsifiable criterion: O1 establishes the reconstructor, while O2 and O3 quantify how its accuracy is set by sensor count, placement, and noise.

\paragraph{O1: An accurate and fast sparse-sensor reconstructor.}
PI-CON must reconstruct the field at engineering grade (defined throughout as a full-band kinetic-energy relative error within $10\%$) from field-reproducible signals only (sensor measurements and PDE residuals, with no DNS full field) and answer arbitrary space--time queries without re-solving the governing equations, using a continuous-time sensor encoder that avoids the sequential time-window training that stabilized MLP-PINNs need on the same high-$Re$ case. On 2D Kolmogorov at $Re = 10\,000$ with $K = 100$ sensors, the criterion is: (i) the full-band KE relative error stays within this target over five seeds, dominated by the energy-rich low-wavenumber band $k \le \lfloor \sqrt{K/\pi} \rfloor = 5$; (ii) a matched $2\times2$ ablation isolates a dominant architectural lever whose main effect reduces KE error by at least two percentage points at $p<0.01$; and (iii) once trained, the model answers any space--time query $(\mathbf{x},t)$ in a single forward pass (without re-integrating the governing equations from an initial condition) and reconstructs a full trajectory at least $5\times$ faster than forward-solving the reference simulation. Fitting is offline and per-case: high-Reynolds chaos ties each realisation to its conditions, so a single shared operator would blur them; cross-case amortisation is not pursued here.

\paragraph{O2: Sensor count sets the achievable resolution.}
The reconstructable wavenumber band must be governed by the sensor count rather than the architecture: $K$ point sensors set a two-dimensional Nyquist-like scale $k_{\max}^{\rm sensor} = \sqrt{K/\pi}$ for the resolvable band, with the effective cutoff limited by sensor information and conditioning rather than architectural capacity. The criterion is: (i) the empirical cutoff $k_{\rm eff}$ lies within $\pm 1.0$ of $k_{\max}^{\rm sensor} \approx 5.64$ at $K = 100$, with mid-band ($5 < k \le 16$) and high-band ($k > 16$) errors reaching the compressed-sensing ceiling of~\eqref{eq:cs_bound} within $[90\%, 110\%]$ and $[95\%, 105\%]$ respectively, indicating that residual error tracks sensor capacity rather than architectural shortfall; and (ii) sensor scaling $K \in \{100, 200, 400\}$ shifts $k_{\rm eff}$ in agreement with the $\sqrt{K/\pi}$ trend.

\paragraph{O3: Placement and noise set the reliability, not the feasibility.}
Reconstruction quality must depend on how sensors are placed and how noisy they are, while staying engineering-grade across reasonable choices. The criterion is: (i) the three placement strategies (a DNS-pivot oracle, a DNS-free LES-derived placement, and random placement) all remain engineering-grade, with the placement-induced standard deviation exceeding the training-seed standard deviation by $\ge 3\times$; and (ii) additive measurement noise up to $10\%$ of the sensor standard deviation keeps the reconstruction engineering-grade in the noise-stress setting.

\paragraph{Contribution.}
PI-CON extends physics-informed operator learning from the dense forward problems where it was developed into the sparse-sensor inverse regime, with each component closing one obstacle from \S\ref{sec:gap}: a closed-form continuous-time (CfC) branch reads each sensor's time history and tolerates irregular measurement clocks (Gaps 2--3); a cross-attention readout with an isotropic distance bias spreads the sparse information across space to return a dense $(u,v,p)$ field at any query $(\mathbf{x},t)$ (Gap 2); and an Augmented-Lagrangian-style adaptive penalty drives the aggregated continuity residual toward zero, keeping the field near-incompressible under sensor-only training (Gap 1). The work also delivers a systematic map of how the sensing configuration governs the result (Gap 4): the sensor count through a Nyquist information ceiling and a band-decomposed error budget, the placement through an oracle/deployable/random comparison, and the measurement noise through a robustness sweep. Among the surveyed methods, the proposed framework is the only one to combine query-anywhere continuous-time evaluation with sensor-only-with-physics training at $Re = 10\,000$: prior sensor-only-with-physics reconstructors are grid-based and demonstrated at $Re \lesssim 10^3$, while the query-anywhere architectures require dense full-field supervision. It attains (O1) while quantifying the (O2)--(O3) sensing limits, with the same architecture remaining feasible two orders of magnitude higher in Reynolds number ($Re = 10^4$--$10^6$, a single-seed feasibility case under retuned sensing and capacity).
